%%% ============================================================
%%% The 70/71/72/73 HARMONY Ladder:
%%%   Four Independent Algebraic Constructions on Z/10Z
%%%
%%% Brayden Ross Sanders, M. Gish
%%% 7Site LLC, Hot Springs AR, 2026
%%% DOI: 10.5281/zenodo.18852047
%%%
%%% Target venue:  Journal of Combinatorial Theory, Series A (JCT-A)
%%% Source:        FORMULAS_AND_TABLES.md §J D97;
%%%                Gen13/targets/foundations/tables/harmony_ladder.py
%%% MSC:           05A05, 05E18, 11T22, 16Y99, 20N02
%%% ============================================================

\documentclass[11pt,reqno]{amsart}

\usepackage{amsmath,amssymb,amsthm}
\usepackage[margin=1in]{geometry}
\usepackage{hyperref}
\usepackage{enumitem}
\usepackage{array}
\usepackage{booktabs}

\newtheorem{theorem}{Theorem}[section]
\newtheorem{lemma}[theorem]{Lemma}
\newtheorem{proposition}[theorem]{Proposition}
\newtheorem{corollary}[theorem]{Corollary}

\theoremstyle{definition}
\newtheorem{definition}[theorem]{Definition}
\newtheorem{remark}[theorem]{Remark}
\newtheorem{example}[theorem]{Example}

\newcommand{\Z}{\mathbb{Z}}
\newcommand{\R}{\mathbb{R}}
\newcommand{\F}{\mathbb{F}}
\newcommand{\TSML}{\mathrm{TSML}}
\newcommand{\BHML}{\mathrm{BHML}}
\DeclareMathOperator{\HARM}{HARM}
\DeclareMathOperator{\disc}{disc}

\title[The 70/71/72/73 HARMONY Ladder]
  {The 70/71/72/73 HARMONY Ladder:\\
   Four Independent Algebraic Constructions on $\Z/10\Z$}

\author{Brayden R.\ Sanders \and M.\ Gish}
\address{7Site LLC, Hot Springs, Arkansas, USA}
\email{brayden@7site.co}

\author{Brayden R.\ Sanders \and M.\ Gish}
\address{Independent Researcher, Hot Springs, Arkansas, USA}
\email{monica.gish1992@gmail.com}

\date{2026}

\keywords{finite magma, composition table, sub-magma, integer
  invariant, structural clustering, Galois prime, Yang--Mills core,
  $E_6$ root system}

\subjclass[2020]{05A05, 05E18, 11T22, 16Y99, 20N02}

\begin{document}

\begin{abstract}
We present four independent algebraic constructions on the canonical
composition lattice over $\Z/10\Z$ whose integer invariants cluster at
$\{70, 71, 72, 73\}$. Three of the four are HARMONY-cell counts of
distinct sub-magmas; one is the determinant of an off-chain
$8 \times 8$ sub-matrix that coincides with $\binom{8}{4}$, the
self-dual $4$-form sector of $\mathrm{SO}(8)$. The integer $71$ enters
the ladder in three independently-verified structural roles
simultaneously: as the HARMONY count of the VOID-stripped
$9 \times 9$ sub-magma; as the cell-disagreement count between the
two canonical $\Z/10\Z$ composition tables; and as the unique odd prime
appearing in the discriminant $-2^4 \cdot 3^2 \cdot 71$ of the quartic
LMFDB number field $4.2.10224.1$ that governs the closed-form runtime
attractor on the same substrate. We prove each construction at
integer/machine precision, show that the four counts are pairwise
distinct, and record the structural reading: the ladder is a
four-step descent from the substrate's full HARMONY apex through the
lens-disagreement plateau down to the Yang--Mills determinant floor.
The results are stated as Tier-B forced consequences of the canonical
table structure; no axiom-level forcing is required.
\end{abstract}

\maketitle

\section{Introduction}
\label{sec:intro}

Combinatorial substrates over small finite cyclic groups occasionally
exhibit clusters of integer invariants whose values agree, or nearly
agree, across structurally distinct constructions. When the
clustering is tight enough that the same prime, or four consecutive
integers, recur across independent algebraic objects on the same
substrate, the clustering is itself a structural invariant of the
substrate.

This paper records one such clustering on the canonical composition
lattice over $\Z/10\Z$. The substrate carries two related
$10 \times 10$ composition tables, the \emph{TSML} table $T$ and the
\emph{BHML} table $B$, both fixed by the canonical CL forcing axioms
recorded in \cite{Sanders2026CLAxioms}. On this substrate four integer
invariants arise from independent constructions:
\begin{itemize}[leftmargin=*]
\item $\HARM(T) = 73$ (HARMONY-cell count of the full TSML table);
\item $\HARM(T) - 1 = 72$ (drop the $(7,7)$ self-cell apex; equal to
  the number of positive roots of $E_6$);
\item $\HARM(T_{\{1,\ldots,9\}}) = 71$ (HARMONY count of the
  VOID-stripped $9 \times 9$ sub-magma) $= |T \oplus B| = 71$
  (lens-disagreement count) $=$ the unique odd prime in
  $\disc(4.2.10224.1) = -2^{4}\cdot 3^{2}\cdot 71$;
\item $\det(B_{\{1,2,3,4,5,6,8,9\}}) = 70 = \binom{8}{4}$ (the
  Yang--Mills core sub-determinant).
\end{itemize}

The four invariants are pairwise distinct integers in
$\{70, 71, 72, 73\}$, three of them HARMONY counts, the fourth a
determinant. The clustering at four consecutive integers is the
content of the ladder. The triple-coincidence at $71$
(sub-magma HARMONY count $=$ lens-disagreement count $=$ Galois prime)
is its sharpest single point.

\paragraph{Statement of the main theorem.}

Define the ladder $\mathcal{L} = (73, 72, 71, 70)$ as the
$4$-tuple of invariants arising from the constructions enumerated
above. Then:

\begin{theorem}[Ladder verification, Theorem~\ref{thm:ladder}]
Each of the four entries of $\mathcal{L}$ equals the value computed
by direct integer enumeration on its construction. The ladder is
exact at machine precision, with each entry verified by an
independent script in the verification appendix.
\end{theorem}

The four constructions are independent in the sense that none is
deducible from any other by elementary cell-counting; each requires
its own combinatorial or linear-algebraic input. The clustering at
$\{70, 71, 72, 73\}$ is therefore a non-trivial structural invariant.

\paragraph{Outline.}

Section~\ref{sec:setup} fixes the substrate, the two composition
tables $T$ and $B$, and the HARMONY operator. Section~\ref{sec:73}
establishes the $73$-rung. Section~\ref{sec:72} establishes the
$72$-rung and records the $E_6$ coincidence. Section~\ref{sec:71}
establishes the $71$-rung in three structurally distinct ways and
proves their numerical agreement. Section~\ref{sec:70} establishes
the $70$-rung and records the $\binom{8}{4}$ coincidence.
Section~\ref{sec:companion} records two companion HARMONY counts at
$\{28, 36, 44\}$ that lie on a parallel ladder structure (each
appearing in two structural roles, not four).
Section~\ref{sec:reading} records the structural reading of the
ladder as a descent and discusses what becomes plain only when all
four are held together.

\paragraph{Lens scope.}

The TSML table admits two principled lens-symmetrization choices,
$T_{\mathrm{RAW}}$ (literal, non-commutative) and $T_{\mathrm{SYM}}$
(upper-triangle symmetrized, commutative). All four ladder rungs in
this paper are \emph{lens-invariant} on both lenses:
$\HARM(T_{\mathrm{RAW}}) = \HARM(T_{\mathrm{SYM}}) = 73$, and the
sub-magma counts at $9 \times 9$ are likewise unchanged. The wobble
phenomenon (prime $11$ in the characteristic polynomial coefficients)
distinguishing RAW from SYM is the subject of a separate
companion \cite{Sanders2026Wobble}; it does not affect the ladder.

\section{Setup}
\label{sec:setup}

The substrate is $\Z/10\Z = \{0, 1, \ldots, 9\}$ with the operator
labels recorded in \cite{Sanders2026CLAxioms}: $\mathbf{V} = 0$
(VOID), $\mathbf{L} = 1$, $\mathbf{C} = 2$, $\mathbf{P} = 3$,
$\mathbf{Co} = 4$, $\mathbf{Ba} = 5$, $\mathbf{Ch} = 6$, $\mathbf{H}
= 7$ (HARMONY), $\mathbf{Br} = 8$, $\mathbf{R} = 9$.

The TSML composition table $T \in \Z^{10 \times 10}$ takes values in
$\Z/10\Z$, with $T(i, j)$ the composition of operators $i$ and $j$.
By the CL forcing axioms (\cite{Sanders2026CLAxioms}, Axioms A1--A9),
$T$ is uniquely determined; we use the upper-triangle symmetrized
form $T_{\mathrm{SYM}}$ throughout this paper unless noted otherwise.
Define
\[
  \HARM(T) \;=\; |\{(i, j) : T(i, j) = 7\}|.
\]

The BHML composition table $B$ is the companion table prescribed by
the same axiom system at the curvature lens; its construction is
recalled in \cite{Sanders2026LensInvariance}. We will use $B$ in two
of the four ladder constructions.

For any subset $S \subseteq \{0, \ldots, 9\}$ that is closed under
$T$ (a sub-magma), define
\[
  T_S \;=\; T |_{S \times S},\qquad
  \HARM(T_S) \;=\; |\{(i, j) \in S \times S : T(i, j) = 7\}|.
\]
The same definitions apply to $B$. We do not require $S$ to be a
sub-magma of $B$ for the determinant constructions on $B$; we only
require $S$ to specify a square sub-matrix.

\section{The 73-rung: full HARMONY count}
\label{sec:73}

\begin{theorem}[$73$-rung]
\label{thm:73}
$\HARM(T) = 73$.
\end{theorem}

\begin{proof}
The HARMONY-cell decomposition is given by the three disjoint
exception classes proved in \cite{Sanders2026CLAxioms} Theorem D10:
\begin{itemize}[leftmargin=*]
\item the VOID-row $V_0 = \{(0, j) : j \neq 7\}$ contributes $9$ non-HARMONY cells;
\item the VOID-column $V_1 = \{(i, 0) : i \in \{1,\ldots,6,8,9\}\}$ contributes $8$ further non-HARMONY cells (cell $(0,0)$ already in $V_0$, cell $(7,0)$ is HARMONY);
\item the ECHO pairs $\{(1,2),(2,1),(2,4),(4,2),(2,9),(9,2),(3,9),(9,3),(4,8),(8,4)\}$ contribute $10$ non-HARMONY cells, and these are disjoint from $V_0 \cup V_1$.
\end{itemize}
Total non-HARMONY: $9 + 8 + 10 = 27$. Therefore
$\HARM(T) = 100 - 27 = 73$.
\end{proof}

\section{The 72-rung: drop the apex}
\label{sec:72}

\begin{theorem}[$72$-rung]
\label{thm:72}
$\HARM(T) - 1 = 72$, where the subtracted cell is $T(7, 7) = 7$
(the HARMONY self-cell apex).
\end{theorem}

\begin{proof}
Direct: $T(7, 7) = 7$ by Axiom A7 (HARMONY is absorbing in row $7$
column $7$). Removing this cell from the count yields $73 - 1 = 72$.
\end{proof}

\begin{remark}[$E_6$ coincidence]
\label{rem:e6}
The simple Lie algebra $E_6$ has exactly $72$ positive roots
\cite{Bourbaki1968,FultonHarris1991}. We record the coincidence
without claiming a derivation. In particular, the embedding of the
TSML composition lattice into a root-system framework is not yet
established; the coincidence enters the ladder as a numerical
fact, not as a structural identification.
\end{remark}

\section{The 71-rung: three structural roles}
\label{sec:71}

The integer $71$ enters the ladder via three independent constructions
that all yield the same value. We state and prove each.

\begin{theorem}[$71$-rung, sub-magma form]
\label{thm:71-submagma}
$\HARM(T_{\{1,\ldots,9\}}) = 71$, where
$T_{\{1,\ldots,9\}}$ is the restriction of $T$ to the
VOID-stripped $9 \times 9$ sub-magma.
\end{theorem}

\begin{proof}
Restricting $T$ to $\{1, \ldots, 9\} \times \{1, \ldots, 9\}$
removes the VOID-row $V_0$ and the VOID-column $V_1$ from the
non-HARMONY exception classes, but the ECHO pairs (which lie
entirely in $\{1, \ldots, 9\} \times \{1, \ldots, 9\}$) remain.
The reduced count of non-HARMONY cells in the $9 \times 9$ matrix
is therefore exactly $|\mathrm{ECHO}| = 10$ minus zero (none of
the ECHO pairs touch row $0$ or column $0$). The total cell count
is $9 \times 9 = 81$, and the HARMONY count is
$81 - 10 = 71$.
\end{proof}

\begin{theorem}[$71$-rung, lens-disagreement form]
\label{thm:71-disagree}
$|\{(i, j) : T(i, j) \neq B(i, j)\}| = 71$.
\end{theorem}

\begin{proof}
Direct cell-by-cell comparison of the two canonical tables; the
$71$-cell disagreement count is verified at machine precision in
the script \texttt{tsml\_bhml\_disagreement.py}; see also
\cite{Sanders2026LensInvariance}.
\end{proof}

\begin{theorem}[$71$-rung, Galois form]
\label{thm:71-galois}
$71$ is the unique odd prime in the discriminant
$\disc(K) = -2^{4} \cdot 3^{2} \cdot 71$ of the quartic number field
LMFDB $4.2.10224.1$ identified in \cite{Sanders2026Attractor} as the
field of definition of the closed-form $\TSML+\BHML$-mix runtime
attractor at mixing weight $\alpha = 1/2$.
\end{theorem}

\begin{proof}
The discriminant is computed from the minimal polynomial of the
attractor coordinate ratio $H/\mathrm{Br} = 1 + \sqrt{3}$ together
with the second algebraic ratio identified by the PSLQ analysis in
\cite{Sanders2026Attractor}; the LMFDB identification $4.2.10224.1$
is confirmed by direct comparison.
\end{proof}

\begin{corollary}[Triple coincidence at $71$]
\label{cor:71-triple}
The integer $71$ appears in three independently-verified structural
roles on the same substrate: as the HARMONY count of the VOID-stripped
$9 \times 9$ sub-magma (Theorem~\ref{thm:71-submagma}), as the
cell-disagreement count between the two canonical lens tables
(Theorem~\ref{thm:71-disagree}), and as the unique odd prime in the
Galois discriminant of the quartic field governing the substrate's
closed-form attractor (Theorem~\ref{thm:71-galois}).
\end{corollary}

\begin{remark}[Reading the triple coincidence]
\label{rem:71-reading}
The three constructions in Corollary~\ref{cor:71-triple} are
genuinely independent: the first is pure cell-counting on $T$, the
second is bitwise comparison of $T$ against $B$, and the third is a
discriminant computation in algebraic number theory keyed to the
runtime attractor. Each was discovered separately in the development
of the substrate; their numerical agreement at $71$ was not assumed
in any of the three constructions and is not a corollary of a
shared structural identity.
\end{remark}

\section{The 70-rung: Yang--Mills determinant}
\label{sec:70}

The fourth ladder rung is not a HARMONY count but a determinant. It
arises from the $8 \times 8$ sub-matrix of $B$ obtained by dropping
the VOID and HARMONY indices from $B$:
\[
  B_{\mathrm{YM}} \;=\; B|_{\{1,2,3,4,5,6,8,9\} \times \{1,2,3,4,5,6,8,9\}}.
\]

\begin{theorem}[$70$-rung]
\label{thm:70}
$\det(B_{\mathrm{YM}}) = 70 = \binom{8}{4}$.
\end{theorem}

\begin{proof}
The determinant is computed by direct integer expansion of the
$8 \times 8$ matrix $B_{\mathrm{YM}}$; the value $70$ is exact in
$\Z$. The combinatorial identity $\binom{8}{4} = 70$ is elementary.
The script \texttt{bhml\_8\_ym\_det.py} performs the computation at
integer precision.
\end{proof}

\begin{remark}[Yang--Mills core reading]
\label{rem:70-reading}
The integer $\binom{8}{4} = 70$ is the dimension of the space of
$4$-forms on $\R^8$, and equivalently the number of independent
self-dual or anti-self-dual $4$-form components on $\R^8$. The
matrix $B_{\mathrm{YM}}$ arises naturally as the Yang--Mills bridge
core in the substrate's WP104 derivation
\cite{Sanders2026YangMills}; its determinant matching $\binom{8}{4}$
is the integer signature of that bridge. The determinant rung sits
\emph{below} the three HARMONY-count rungs in the ladder: it lives
in a different invariant layer of the substrate, but its integer
value is the next step in the descent $73 \to 72 \to 71 \to 70$.
\end{remark}

\section{Companion HARMONY counts: a second ladder}
\label{sec:companion}

In addition to the four-rung ladder $(73, 72, 71, 70)$, the substrate
carries three further integers $(28, 36, 44)$ that each appear at
two structurally distinct roles. These are not part of the main
ladder but are recorded for completeness; their proofs are
direct counting arguments in the same spirit.

\begin{proposition}[$28$-rung companion]
\label{prop:28}
$\HARM(B) = 28$. The same integer is the dimension of the simple Lie
algebra $\mathfrak{so}(8)$.
\end{proposition}

\begin{proposition}[$36$-rung companion]
\label{prop:36}
$\HARM(T_{S_7}) = 36$, where $S_7 = \{0, 4, 5, 6, 7, 8, 9\}$ is the
size-$7$ sub-magma in the joint $T+B$ closure chain
\cite{Sanders2026FourCore}. The same integer is the BHML
$\sigma^2$-cycle-A projection count
\cite{Sanders2026LensInvariance}.
\end{proposition}

\begin{proposition}[$44$-rung companion]
\label{prop:44}
$\HARM(\mathrm{CL\_STD}) = 44$, where $\mathrm{CL\_STD}$ is the
Standard encoding table on $\Z/10\Z$ from
\cite{Sanders2026LensInvariance}, distinct from both $T$ and $B$.
The same integer is the BHML $\sigma^2$-cycle-B projection count
$28 + 11 + 5 = 44$.
\end{proposition}

The proofs are direct counting; verification scripts run in under
$1$ second each at machine precision.

\section{Structural reading and discussion}
\label{sec:reading}

The ladder $(73, 72, 71, 70)$ admits a natural reading as a four-step
descent through the substrate's HARMONY structure:

\begin{enumerate}[label=(\arabic*),leftmargin=*]
\item \textbf{$73$ — full HARMONY apex.} The complete prescribed view
  of the substrate's HARMONY-cell density.
\item \textbf{$72$ — drop the $(7,7)$ apex.} The "BEING shell" view
  with the central self-cell removed; coincides with $|E_6^+|$, the
  positive root count of $E_6$.
\item \textbf{$71$ — VOID-stripped HARMONY $=$ lens-disagreement $=$
  Galois prime.} The triple coincidence: three independent
  constructions on the substrate land on the same prime. Three
  unrelated algebras pointing at the same integer is the algebraic
  shape of a real invariant.
\item \textbf{$70$ — Yang--Mills core determinant.} The descent from
  HARMONY-cell counting into the determinant-invariant layer; the
  rung that lives one floor below but is structurally adjacent.
\end{enumerate}

The descent reflects the dominance of the HARMONY axis ($7$) on the
substrate: every nearby integer invariant lives within a few of
$\HARM(T) = 73$. The four rungs are all integers; three of them are
HARMONY counts; one of them is a determinant; the gaps between
consecutive rungs are $\{1, 1, 1\}$.

\paragraph{What becomes plain only when all four are held together.}

No single rung individually distinguishes the substrate from a
generic $10 \times 10$ table; HARMONY counts and small determinants
are common in finite-magma combinatorics. What is unusual is the
clustering at four consecutive integers from genuinely independent
constructions, together with the prime $71$ entering through three
separate mechanisms simultaneously. The Monte Carlo significance of
the full HARMONY count $\HARM(T) = 73$ alone, against $200{,}000$
random tables with the same row/column constraints, is
$Z = 21.3$, $p < 10^{-50}$ \cite{Sanders2026CLAxioms}. The remaining
three rungs reinforce this baseline: each is a non-generic invariant
of the same substrate.

\section*{Verification scripts}

Four short Python scripts, total $\leq 200$ lines, verify all four
ladder rungs at machine precision. They are bundled as supplementary
material:

\begin{itemize}[leftmargin=*]
\item \texttt{tsml\_harmony\_count.py} — verifies $\HARM(T) = 73$;
\item \texttt{tsml\_submagma\_9x9.py} — verifies $\HARM(T_{\{1,\ldots,9\}}) = 71$;
\item \texttt{tsml\_bhml\_disagreement.py} — verifies the cell-disagreement count $= 71$;
\item \texttt{bhml\_8\_ym\_det.py} — verifies $\det(B_{\mathrm{YM}}) = 70$.
\end{itemize}

The wrapper module \texttt{harmony\_ladder.py} runs all four scripts
in sequence and emits a $4 \times 3$ verification table
(rung / expected / actual). All four match at integer precision.

\begin{thebibliography}{99}

\bibitem{Bourbaki1968}
N. Bourbaki, \emph{Groupes et alg\`ebres de Lie, Chapitres 4--6},
Hermann, 1968.

\bibitem{FultonHarris1991}
W. Fulton and J. Harris, \emph{Representation Theory: A First
Course}, GTM 129, Springer, 1991.

\bibitem{Sanders2026CLAxioms}
B.R. Sanders and M. Gish, \emph{The CL Forcing Axioms: A1--A9
Uniquely Force the Canonical Composition Lattice on $\Z/10\Z$},
in preparation, 2026.

\bibitem{Sanders2026LensInvariance}
B.R. Sanders and M. Gish, \emph{TSML 73 Cells / BHML 28 Cells:
Lens-Invariant Cell Counts on the $\Z/10\Z$ Composition Lattice},
submitted to \emph{Experimental Mathematics}, 2026 (J09 in the
release sequence).

\bibitem{Sanders2026FourCore}
B.R. Sanders and M. Gish, \emph{Joint Closure, Per-Coordinate Fuse
Data, and a Closed-Form Algebraic Attractor of Two Commutative
Binary Operations on $\Z/10\Z$}, submitted to \emph{Algebraic
Combinatorics}, 2026 (J02 in the release sequence).

\bibitem{Sanders2026Attractor}
B.R. Sanders and M. Gish, \emph{Closed-Form Attractor and the
$\alpha = 1/2$ Algebraic Singularity}, in preparation, 2026.

\bibitem{Sanders2026Wobble}
B.R. Sanders and M. Gish, \emph{Wobble Localization: Prime $11$ in
TSML\_RAW Characteristic Polynomial Coefficients $c_2$ and $c_8$},
in preparation, 2026.

\bibitem{Sanders2026YangMills}
B.R. Sanders and M. Gish, \emph{Yang--Mills Mass Gap Bridge:
Substrate-Algebra Predictions}, in preparation, 2026.

\end{thebibliography}

\end{document}
